\chapter{Pendahuluan}

\section{Latar Belakang}
Volatilitas merupakan masukan utama dalam manajemen risiko, penetapan harga
derivatif, dan alokasi aset, sehingga ketepatan peramalannya berdampak langsung
pada keputusan finansial. Volatilitas tidak dapat diamati secara langsung,
tetapi dapat diukur melalui \textit{realized variance} (RV), yaitu jumlah
kuadrat \textit{return} intraday dalam satu hari. Andersen dan Bollerslev
menunjukkan bahwa proksi volatilitas yang berisik membuat model tampak buruk,
sedangkan RV yang dihitung dari data intraday memberikan ukuran yang jauh lebih
akurat \cite{andersen1998}. Di antara model peramalan RV, model
\textit{Heterogeneous Autoregressive} (HAR-RV) yang diusulkan Corsi menjadi
pembanding yang sulit dikalahkan karena strukturnya sederhana namun mampu
menangkap sifat memori panjang volatilitas melalui rata-rata RV harian,
mingguan, dan bulanan \cite{corsi2009}.

Bitcoin, aset kripto pertama yang diperkenalkan melalui sistem pembayaran
elektronik \textit{peer-to-peer} \cite{nakamoto2008}, merupakan objek yang
menarik untuk peramalan volatilitas karena karakteristik pasarnya berubah
substansial dalam satu dekade terakhir. Do membagi perkembangan pasar Bitcoin
periode 2012--2025 ke dalam tiga era, yaitu \textit{niche},
\textit{transition}, dan \textit{institutional}, dengan persistensi RV yang
berbeda antarera dan respons \textit{illiquidity} terhadap guncangan
volatilitas yang lebih kecil serta lebih singkat pada era \textit{institutional}
\cite{do2026}. Perubahan tersebut tidak berlangsung mulus. Efisiensi pasar
Bitcoin berganti-ganti dari waktu ke waktu, dan koefisien likuiditas dapat
berbalik tanda sebelum dan sesudah pandemi COVID-19 \cite{mokni2024}. Kondisi
ini memunculkan kemungkinan bahwa hubungan antara informasi pasar dan
volatilitas yang dipelajari sebuah model pada satu periode tidak lagi berlaku
pada periode berikutnya.

Perubahan karakteristik pasar perlu dibedakan dari klaim bahwa volatilitas
Bitcoin tidak dapat diramal. Literatur efisiensi pasar Bitcoin menguji apakah
\textit{return} dapat diprediksi dari informasi masa lalu
\cite{yi2023,mokni2024,do2026}. Pasar yang mendekati efisien dalam pengertian
tersebut tetap dapat memiliki volatilitas yang terprediksi, karena volatilitas
cenderung mengelompok: hari yang bergejolak cenderung diikuti hari yang
bergejolak. Sifat inilah yang dimanfaatkan model HAR-RV \cite{corsi2009}.
Dengan demikian, peramalan volatilitas Bitcoin tetap bermakna meskipun
efisiensi \textit{return}-nya meningkat.

Berbagai penelitian telah membandingkan model \textit{machine learning} (ML)
dengan HAR-RV dalam meramalkan volatilitas. Christensen, Siggaard, dan Veliyev
menjadi salah satu rujukan utama perbandingan ML dan HAR \cite{christensen2023}.
Pada Bitcoin, Huang, Sangiorgi, dan Urquhart melaporkan bahwa LSTM hanya unggul
tipis atas HAR, sementara HAR lebih baik ketika terjadi lonjakan volatilitas
\cite{huang2024}. Dudek dkk. membandingkan sejumlah model pada periode
2019--2021 dan melaporkan kinerjanya sebagai satu angka agregat; kinerja pada
hari dengan RV terendah dan tertinggi hanya dibedakan secara deskriptif tanpa
uji formal \cite{dudek2024}. Pola serupa tampak pada penelitian lain: hampir
semuanya memakai satu periode uji dan melaporkan hasil agregat, dan sebagian
memakai fungsi \textit{loss} yang tidak robust terhadap proksi volatilitas yang
berisik \cite{huang2024,wang2023,dudek2024,oprea2026,patton2011}.

% CEK: Christensen dkk. (2023) baru metadata. Klaim dibatasi pada "rujukan
% perbandingan ML dan HAR"; perinci setelah full text dibaca.

Pelaporan hasil agregat menyimpan persoalan. Bahwa kinerja relatif model dapat
berubah antarperiode uji sudah dikenal dalam literatur peramalan
\cite{tashman2000}, dan literatur ML menegaskan bahwa rata-rata kinerja dapat
menyembunyikan perilaku model sepanjang waktu \cite{gama2014}. Literatur
\textit{forecasting under instability} menyediakan alat uji untuk mendeteksi
perubahan kinerja relatif tersebut \cite{giacomini2010,rossi2013,martins2016}.
Alat uji ini sudah diterapkan pada aset kripto oleh Trucíos dan Taylor, tetapi
untuk peramalan \textit{Value at Risk} dan \textit{Expected Shortfall} dengan
model ekonometrik, bukan untuk perbandingan ML dan HAR pada RV
\cite{trucios2023}. Dengan kata lain, belum jelas apakah keunggulan atau
kekalahan ML terhadap HAR pada volatilitas Bitcoin bersifat stabil, atau muncul
dan hilang mengikuti kondisi pasar.

% CEK: Giacomini & Rossi (2010) baru metadata; Trucíos & Taylor (2023) baru
% abstrak.

Persoalan berikutnya adalah cara model diperbarui, yaitu seberapa banyak data
lama yang masih layak dipakai untuk melatih ulang model. Pilihan antara
\textit{rolling window} dan \textit{expanding window} berpengaruh terhadap
akurasi peramalan volatilitas \cite{feng2024window}, termasuk pada aset kripto
\cite{akgun2025}. Namun, Akgun dan Gulay melaporkan perbandingan kedua skema
tersebut secara rinci hanya untuk model tipe GARCH, tanpa HAR sebagai
pembanding dan tanpa uji signifikansi \cite{akgun2025}, sedangkan Feng, Qi, dan
Lucey memvariasikan \textit{hyperparameter} model ML, bukan panjang data latih
\cite{feng2024ml}.
Chassot dan Audrino menunjukkan pada 1.445 saham bahwa model ML gagal
mengalahkan HAR yang diberi skema \textit{fitting} yang tepat, sehingga
perbandingan ML dengan HAR dapat bias apabila HAR tidak diperlakukan setara
\cite{chassot2026}. Belum jelas apakah strategi pembaruan memengaruhi kinerja
relatif ML terhadap HAR secara berbeda pada era pasar yang berbeda.

% CEK: Feng (2024) baru metadata; Feng, Qi & Lucey (2024) dan Chassot & Audrino
% (2026) baru abstrak. Klaim pembeda di paragraf ini bergantung pada full text
% keduanya (checklist novelty).

Penelitian ini juga memperhatikan persoalan pengukuran. \textit{Loss} absolut
otomatis meningkat ketika volatilitas tinggi karena skala targetnya membesar,
sehingga kenaikan \textit{loss} tidak dapat langsung ditafsirkan sebagai
penurunan kualitas model. Untuk itu, objek analisis penelitian ini adalah
selisih \textit{loss} harian antara model ML dan HAR-RV pada hari yang sama,
yang menetralkan efek skala tersebut. Fungsi \textit{loss} utama yang digunakan
adalah QLIKE karena robust terhadap proksi volatilitas yang berisik dan tidak
bergantung pada skala \cite{patton2011}.

Berdasarkan uraian tersebut, penelitian ini menguji stabilitas kinerja relatif
LightGBM \cite{ke2017} terhadap HAR-RV dalam meramalkan \textit{realized
variance} harian BTC/USDT spot di Binance yang dihitung dari \textit{return}
5 menit. Penelitian menggunakan desain $2 \times 2$ (HAR, HAR-X, LGBM-HAR, dan
LGBM-X) untuk memisahkan kontribusi non-linearitas dari kontribusi informasi
tambahan, serta membandingkan strategi pembaruan \textit{static},
\textit{expanding}, dan \textit{rolling} dengan perlakuan yang setara bagi
HAR-RV. Penelitian ini memperluas Dudek dkk. \cite{dudek2024} dengan menguji
kestabilan kinerja relatif sepanjang waktu, faktor pasar yang berasosiasi
dengan perubahannya, dan pengaruh strategi pembaruan model. Kontribusi
penelitian terletak pada rancangan evaluasi model ML pada data yang berubah,
yang memadukan strategi pembaruan dari literatur \textit{concept drift} dengan
uji kestabilan dari ekonometrika; kedua literatur ini dicatat masih kurang
terhubung \cite{suarez2023}. Kontribusi tersebut bukan algoritma baru maupun
teori keuangan.

\section{Rumusan Masalah}
Berdasarkan latar belakang di atas, rumusan masalah penelitian ini adalah
sebagai berikut.
\begin{enumerate}
  \item Sejauh mana kinerja relatif model \textit{machine learning} terhadap
        HAR-RV dalam meramalkan \textit{realized variance} Bitcoin stabil
        sepanjang periode pengujian?
  \item Apakah variasi kinerja relatif tersebut berasosiasi dengan
        karakteristik pasar dan pergeseran distribusi fitur?
  \item Apakah strategi pembaruan model dan panjang \textit{window} data latih
        memengaruhi kinerja relatif \textit{machine learning} terhadap HAR-RV
        secara berbeda pada era pasar yang berbeda?
\end{enumerate}

% CEK: Dokumen riset menulis "keunggulan relatif"; di sini diselaraskan menjadi
% "kinerja relatif" mengikuti istilah judul v2.5. Substansi RQ tidak berubah.
% Konfirmasi ke pembimbing.

\section{Tujuan Penelitian}
Tujuan penelitian ini adalah sebagai berikut.
\begin{enumerate}
  \item Menguji secara formal kestabilan kinerja relatif model \textit{machine
        learning} terhadap HAR-RV, dengan dekomposisi $2 \times 2$ untuk
        melihat sumber perbedaan kinerja, yaitu non-linearitas atau informasi
        tambahan.
  \item Mengidentifikasi kondisi pasar, yang diukur melalui karakteristik pasar
        dan pergeseran distribusi fitur, ketika model \textit{machine learning}
        lebih unggul atau lebih buruk daripada HAR-RV.
  \item Mengevaluasi apakah data dari era awal membantu atau merugikan kinerja
        relatif \textit{machine learning} terhadap HAR-RV, dengan perlakuan
        \textit{window} yang setara bagi kedua model.
\end{enumerate}

\section{Manfaat Penelitian}

\subsection{Manfaat Teoretis}
Penelitian ini memberikan bukti empiris mengenai kestabilan kinerja relatif
model \textit{machine learning} terhadap HAR-RV pada peramalan volatilitas
Bitcoin lintas era pasar, beserta sumber perbedaan kinerjanya melalui
dekomposisi $2 \times 2$. Penelitian ini juga menyediakan rancangan evaluasi
temporal, meliputi \textit{pipeline} data, strategi pembaruan model, dan uji
kestabilan, yang dapat direplikasi pada aset dan model lain.

\subsection{Manfaat Praktis}
Bagi praktisi dan pengembang sistem peramalan, hasil penelitian ini membantu
menentukan apakah model \textit{machine learning} layak digunakan dibandingkan
model sederhana, dalam kondisi pasar apa kinerjanya cenderung berubah, dan
seberapa banyak data historis yang sebaiknya dipakai ketika model dilatih
ulang.

\section{Batasan Masalah}

% CEK: Sub-bab ini tidak ada di template prodi. Pertahankan hanya bila
% pembimbing menyetujui; bila tidak, pindahkan isinya ke Bab III (ruang lingkup
% penelitian).

Agar penelitian terarah, batasan masalah ditetapkan sebagai berikut.
\begin{enumerate}
  \item Data yang digunakan adalah harga BTC/USDT spot dari arsip publik
        Binance dengan resolusi 5 menit, sejak awal ketersediaan data pada
        Agustus 2017 hingga September 2026. Kesimpulan tidak digeneralisasi ke
        pasar \textit{perpetual} maupun bursa lain.
  \item Target peramalan adalah \textit{realized variance} harian untuk horizon
        satu hari ke depan.
  \item Model yang dibandingkan adalah HAR-RV, HAR-X, LGBM-HAR, dan LGBM-X,
        dengan \textit{naive persistence} sebagai pembanding dasar. Model
        \textit{deep learning} tidak termasuk dalam cakupan penelitian.
  \item Periode pengujian mencakup era \textit{transition} dan
        \textit{institutional}; era \textit{niche} tidak tercakup karena data
        spot Binance baru tersedia sejak 2017.
  \item Seluruh variabel pasar diturunkan dari data OHLCV. Penelitian bersifat
        asosiatif sehingga tidak mengklaim hubungan kausal, tidak mengklaim
        pembuktian \textit{concept drift} atau \textit{covariate shift} dalam
        arti formal, dan tidak menjelaskan struktur mikro pasar Bitcoin.
\end{enumerate}
